\section{Introduction}
\label{sec:intro}

In a review of FDA records for twelve antidepressants, 94\% of the trials in the
published literature were positive, against 51\% in the FDA's own analysis, and
the published effect sizes were 32\% larger overall \citep{turner2008selective}.
The pattern is not confined to medicine: in a set of peer-reviewed
social-science experiments, strong results were 60 percentage points more likely
to be written up than null results \citep{franco2014publication}. Selection also
acts inside a study. Comparing trial protocols with their publications, 40--62\%
of trials had at least one primary outcome changed, introduced or omitted, and
statistically significant outcomes were more likely to be fully reported
\citep{dwan2008systematic}. This is publication bias and reporting bias, the
file drawer problem of \citet{rosenthal1979file}, and it needs no fraud: small
studies, flexible analyses and competing interests already make many published
findings likely to be false \citep{ioannidis2005why}, and incentives that reward
publication can select for weak methods without anyone cheating deliberately
\citep{smaldino2016natural}. Authors respond to the credit they receive, and a
literature that credits positive, attention-getting results teaches them to
report only those.

Behind these figures lies a problem of rewarding exploration. Agents that search
on behalf of a group bear the cost of exploring, but what they find can save
everyone else the same work, so a group gains only when its members explore and
then report truthfully what they find. Models of the division of cognitive labor show that
the reward for discovery shapes which problems a community chooses to pursue
\citep{kitcher1990division, strevens2003priority}, and the same dependence
appears when a planner decides what to disclose to agents whose own experiments
create the information others rely on \citep{kremer2014wisdom}. How to reward
such exploration is a credit assignment problem, the problem of crediting each
agent for its own contribution to a shared return \citep{foerster2018coma}. A
growing family of methods now rewards agents directly for their effect on others:
social-influence rewards credit an agent for changing other agents' actions
\citep{jaques2019social}, and learned incentive functions let agents reward one
another to shape behavior \citep{yang2020lio}. Whenever credit flows to whoever
is relied on, it determines not only how much a group explores but whether its
members report what they find honestly.

Theory explains why such failures persist. Conflicts of interest limit how much
information survives strategic communication \citep{crawford1982strategic},
receivers who take messages at face value are deceived by inflated language
\citep{kartik2007credulity}, and a sender who controls what is revealed can
steer a rational receiver's choice \citep{kamenica2011bayesian}. Citation counts
can be inflated outright, as when editors pressure authors to cite their journal
\citep{wilhite2012coercive}, and reward-maximizing learning agents with no
built-in intent to deceive have been observed to distort their signals once
distortion is profitable \citep{alon2023disinformation}.

Intuitively, there are six common designs for the credit. However, each falls
short.

\paragraph{Credit from the shared return.}

Every agent receives the group's return, the default in cooperative multi-agent
learning, where a shared reward does not say which agent contributed to it
\citep{foerster2018coma}. An explorer bears the full cost of exploring, while its
own contribution to a return shared by many agents is hard to discern, so
exploring and reporting are under-supplied. It also cannot tell a true report
from an untruthful one, since both change the shared return only through what
others later do with them.

\paragraph{Credit for keeping exploring.}

Explorers are credited for each exploration, whatever it yields. This sustains exploration, but it is blind to what is reported: a true report, an exaggerated one and silence all earn the same, so honesty has no advantage, and credit is spent on explorations whose findings never reach the group. Untruthful reports then harm the group that relies on them: they inflate apparent effects, as in
the antidepressant trials above \citep {turner2008selective}, and they add to avoidable waste in research \citep{chalmers2009avoidable}.

\paragraph{Credit for only being followed (Popularity credit).}

Explorers are credited for attention, not accuracy. It encourages agents to explore
more as they get rewarded more. A inflated claim attracts followers as readily as a real one, so exaggeration is profitable and a learner finds it. This is the same situation that celebrities may promote low-quality products without careful consideration, prioritizing their own financial gain over the well-being of consumers or reporting bias aforementioned.

\paragraph{Credit for only being accurate.}

Explorers are credited only for the verified gain of those who rely on a claim. This
leaves a true report that nobody uses with nothing, which is the incentive
behind the file drawer, and a community can stall in a state where no one
reports because no one listens. Such credit also never checks an unused claim,
so truth, exaggeration and silence all earn the same. Accuracy-based remedies
from the literature do not close this gap: peer prediction makes honesty an
equilibrium rather than a dominant strategy \citep{miller2005peer}, and proper
scoring rules require the outcome to be observed \citep{gneiting2007strictly}.

\paragraph{A cost for lying alone.}

A fine charged when an audit finds a report false \citep{becker1968crime}. It
helps reduce lying rate of the agents. Nevertheless, it is not enough either: a
lie that is followed still collects the followers' credit, so the fine has to
outbid it, and a true report that nobody follows still earns nothing, so truth
is no better than silence. Moreover, a lie can be profitable outside the credit channel,
as when a hyped claim attracts followers whose attention benefits the claimant,
so the penalty must outweigh a temptation the designer does not control. With
smooth lying costs, \citep{kartik2009StrategicCommunicationLying} shows that
truthful reporting is never an equilibrium, and that inflation shrinks with cost
intensity but persists.

\paragraph{Reputation without credit.}

Receivers stop following senders who were wrong before, and give nothing. This
learns only from claims that were checked, and a claim that nobody follows is
never checked, so unused lies leave a sender's reputation untouched. It also
gives an honest explorer no reason to report: a new sender starts without trust
and may not be followed, so a true report can go unused and earn nothing.

These failures set our goal. We ask for a credit rule under which reporting the
truth is each agent's best action whatever the others do with the report,
including when a lie or a withheld result carries a private gain the designer
neither funds nor observes. Because receivers in learning systems, like
readers of a literature, keep changing how they respond, the guarantee cannot
rest on an equilibrium or on a model of receiver behavior; it must hold for
every receiver rule, so truth must beat the response most favorable to a lie.

Our key idea is to reward only verified value and to make every claim face a
check. We propose Truthful Verified Credit (\tvc{}), which wraps any host credit
scheme with three parts, each rewarding truth at least as much as a lie in every
case. Verified gain credit rewards the realized improvement of those who rely on a
report, capped, and only when the report passes verification. A checked score
bonus rewards every report that is checked, whether by a follower or by an
independent random audit, so a true report, including a negative one, earns
credit even when nobody acts on it. An audit fine is charged when an audit finds
a report false, and we derive the fine needed to outweigh a lie's outside gain.

We study \tvc{} in two settings. In a multi-agent bandit, where agents search
arms of unknown value on behalf of a group, popularity credit drives over
90\% of reports to be exaggerated, while \tvc{} brings it below one percent at
comparable welfare. In a congestion setting built on a multi-agent pathfinding (MAPF)
benchmark \citep{skrynnik2025pogema}, robots can post fake jams to clear their
own routes: self-interested learners lie at 91--99\% of opportunities without
credit, \tvc{} holds the rate near 0.2\%, and throughput rises about threefold.

% We summarize our contributions as follows.
% \begin{itemize}
%   \item We formalize collective exploration with reliance-based credit and show that rewarding attention breeds hype and silence; when information drives congestion, hype becomes a social dilemma in which every agent hyping lowers group output by 65\% (\cref{sec:problem,sec:experiments}).
%   \item We propose \tvc{}, show that truth earns more than any lie for every receiver rule, and show by ablation that each of its parts is needed (\cref{sec:why,sec:experiments}).
%   \item We derive the audit fine required when lies carry outside gains and confirm the predicted threshold empirically (\cref{eq:fine,sec:experiments}).
%   \item We evaluate \tvc{} in a multi-agent bandit and in multi-agent pathfinding, and report its limits: \tvc{} elicits truth rather than usefulness, so its unconditional bonus can encourage reports that carry no value, and it raises welfare only where honest information changes what others do (\cref{sec:experiments,sec:conclusion}).
% \end{itemize}
